\documentclass[11pt,a4paper]{article}

\usepackage[margin=0.78in]{geometry}
\usepackage[T1]{fontenc}
\usepackage[utf8]{inputenc}
\usepackage{lmodern}
\usepackage{microtype}
\usepackage{graphicx}
\usepackage{booktabs}
\usepackage{tabularx}
\usepackage{longtable}
\usepackage{array}
\usepackage{amsmath}
\usepackage{siunitx}
\usepackage{xcolor}
\usepackage{hyperref}
\usepackage{caption}
\usepackage{float}
\usepackage{enumitem}

\hypersetup{
  colorlinks=true,
  linkcolor=blue!50!black,
  urlcolor=blue!50!black,
  citecolor=blue!50!black,
  pdftitle={Task 1: Forecasting Across Heterogeneous Sensors},
  pdfsubject={AI651 Assignment 1}
}
\setlength{\parskip}{0.45em}
\setlength{\parindent}{0pt}
\setlength{\emergencystretch}{2em}
\captionsetup{font=small,labelfont=bf}
\newcommand{\mse}{\ensuremath{\mathrm{MSE}}}
\newcommand{\rmse}{\ensuremath{\mathrm{RMSE}}}
\newcommand{\units}{\ensuremath{\mathrm{m/s^2}}}
\newcommand{\fig}[3]{
  \begin{figure}[H]
    \centering
    \includegraphics[width=\linewidth,height=0.40\textheight,keepaspectratio]{#1}
    \caption{#2}
    \label{#3}
  \end{figure}
}

\title{\textbf{Task 1: Forecasting Across Heterogeneous Sensors}\\
  \large AI651 -- Deep Learning for Space, Time and Graphs\\
  \normalsize Assignment 1}
\author{Student: \rule{2.7in}{0.4pt} \qquad Student ID: \rule{1.3in}{0.4pt}}
\date{Fall 2026}

\begin{document}
\maketitle

\begin{center}
\textbf{Repository URL:} \rule{4.2in}{0.4pt}
\end{center}

\begin{abstract}
This report evaluates context-dependent forecasting for a synthetic four-station production line.
The model receives a 96-sample local history and predicts 48 future samples. Results compare
shared, sensor-specific, and period-routed ridge regression with Raw Attention, moving-average
decomposition, and a delay mixer. On the full-preset, single-seed run, period-routed ridge has the
lowest validation error for both established and held-out populations, while the combined
Autoformer-inspired model is the strongest neural model in aggregate. The report distinguishes
validation-based choices from final test evaluation and documents the limitations of this one
synthetic run.
\end{abstract}

\clearpage
\tableofcontents
\newpage

\section{Task, data, and evaluation protocol}

The simulated production line has four vibration stations sampled at 20\,Hz. Products A, B, and C
have respective nominal period ranges of 14--18, 21--27, and 28--36 samples. The signal combines a
slow baseline/drift, a repeating vibration, and noise. Stations 1--3 are established; station 4 is
held out from model fitting. A forecaster sees only one station's 96-sample context and predicts
the next 48 samples. Windows do not cross operating-regime boundaries.

Training, validation, and final test targets lie in the disjoint intervals
$[0,9600)$, $[9600,12480)$, and $[12480,15360)$, respectively. Contexts are anchored at their
last observed value and scaled using established-station training data only; predictions are
restored to physical units for evaluation. The saved notebook setup records \texttt{PA1\_PRESET=full},
data seed 0, model seed 0, and CPU device. Neural runs report 6,000 optimization steps.

\textbf{Metrics.} RMSE is reported in \units{} and MSE in $(\units)^2$. Lower is better. The
validation set selects models and deployment choices; the final test set is shown only after those
choices are made.

\section{Part 1: context-dependent forecasting}

\subsection{Response 1A: hypothesis, results, and interpretation}

\textbf{Timing note.} The hypothesis below was entered into the notebook after the saved experiment
outputs existed. It is therefore reported transparently as a retrospective hypothesis, not as a
preregistered prediction.

The expected ordering was period-routed ridge, then Raw Attention, then approximately tied shared
and sensor-specific ridge on established stations. The same ordering was expected on station 4,
without sensor-specific ridge. The rationale was that period routing estimates pace from the
observed context and selects a fitted map, whereas fixed ridge maps compromise across changing
paces; Attention must learn context-dependent mixing.

\subsubsection{Output 1.1: validation comparison}

\begin{table}[H]
\centering
\caption{Output 1.1: full-preset validation results. Product entries are RMSE.}
\label{tab:part1}
\small
\begin{tabularx}{\linewidth}{>{\raggedright\arraybackslash}X*{6}{>{\centering\arraybackslash}X}}
\toprule
Model & Est. RMSE & Est. MSE & Held-out RMSE & Held-out MSE & Parameters & Fit (s)\\
\midrule
Shared ridge & 0.767 & 0.588 & 0.876 & 0.767 & 4,608 & 0.002\\
Sensor-specific ridge & 0.768 & 0.590 & unavailable & unavailable & 13,824 & 0.003\\
Period-routed ridge & \textbf{0.166} & \textbf{0.028} & \textbf{0.173} & \textbf{0.030} & 59,904 & 0.006\\
Raw Attention & 0.431 & 0.186 & 0.612 & 0.374 & 368,368 & 760.032\\
\bottomrule
\end{tabularx}
\end{table}

\begin{table}[H]
\centering
\caption{Output 1.1 product-wise validation RMSE by population.}
\label{tab:part1-products}
\small
\begin{tabular}{llrrrr}
\toprule
Population & Model & Product A & Product B & Product C & Units\\
\midrule
Established & Shared ridge & 0.805 & 0.854 & 0.621 & \units\\
Established & Sensor-specific ridge & 0.804 & 0.852 & 0.630 & \units\\
Established & Period-routed ridge & \textbf{0.181} & \textbf{0.164} & \textbf{0.152} & \units\\
Established & Raw Attention & 0.388 & 0.482 & 0.418 & \units\\
Held-out & Shared ridge & 0.909 & 0.982 & 0.713 & \units\\
Held-out & Period-routed ridge & \textbf{0.184} & \textbf{0.182} & \textbf{0.150} & \units\\
Held-out & Raw Attention & 0.647 & 0.580 & 0.607 & \units\\
\bottomrule
\end{tabular}
\end{table}

On established stations, the measured ordering is period-routed ridge $<$ Raw Attention $<$
shared ridge $\approx$ sensor-specific ridge. On station 4 it is period-routed ridge $<$
Raw Attention $<$ shared ridge; sensor-specific ridge is unavailable because no station-4
matrix is fitted. Relative to shared ridge, period routing lowers aggregate RMSE by about 78\%
on established stations and 80\% on station 4. Raw Attention lowers it by about 44\% and 30\%,
respectively.

Across products, routed ridge and Raw Attention retain their relative ordering. Shared and
sensor-specific ridge are near-tied and swap order: sensor-specific is marginally better for A
and B, while shared is better for C. Station-specific fitting does not solve pace changes within
a station. Period routing provides a useful explicit inductive bias by estimating period from
the current context and choosing a pace-band map. Attention also depends on context, but in this
run its flexibility does not beat period routing.

\subsection{Response 1B: why forecasting must depend on context}

\subsubsection{Outputs 1.2 and 1.3: context and Attention}

\fig{extracted_images/01-cell-12-output-02.png}
{Output 1.2: forecasts for Station 1 under Product A and Product C, with different actual periods
(16.135 and 30.093 samples). The fixed ridge curves compromise in phase/shape; the context-routed
forecast follows the continuations more closely.}
{fig:regime-forecasts}

In the two plotted validation windows, shared-ridge window RMSE is 0.779 for Product A and 0.448
for Product C; period-routed ridge gives 0.232 and 0.150, and Raw Attention gives 0.290 and
0.378. These are individual examples, not aggregate performance estimates. The same station faces
different periods, so station identity alone cannot specify the phase advance needed for the next
48 samples. A fixed linear map must compromise across these incompatible continuations.

\fig{extracted_images/02-cell-13-output-02.png}
{Output 1.3: learned Attention weights respond differently to Product A and Product C contexts.
This demonstrates context dependence, not recovery of the physical period.}
{fig:attention}

The learned query, key, and value projections are fixed after training, but their outputs depend
on the input context. Thus the query-key scores and softmax weights, and therefore the value-mixing
operation, change from one context to another. The saved mean absolute weight difference is
0.0113 and the maximum is 0.4250. This does not establish that Attention has recovered the true
period; it shows only that its mixing weights depend on the input.

\section{Part 2: separating slow structure from repetition}

\subsection{Response 2: decomposition, width selection, and evidence}

For odd width $m=2q+1$, the trend is the centered moving average
\[
T_t=\frac{1}{m}\sum_{u=-q}^{q}x_{t+u},\qquad S_t=x_t-T_t.
\]
At boundaries the nearest endpoint is replicated. The streams have the same length and satisfy
$S+T=x$. The trend is sent to a linear forecasting head, while the remainder is embedded and
processed by the encoder; decomposition also follows each residual sublayer.

\fig{extracted_images/03-cell-18-output-02.png}
{Output 2.1: candidate moving-average widths on a Station 2/Product C context. Wider width 41
leaves a stronger oscillatory residual and yields a clear period peak after decomposition.}
{fig:decomposition}

\begin{table}[H]
\centering
\caption{Output 2.1: candidate trend widths and training-only oracle diagnostic.}
\label{tab:widths}
\small
\begin{tabular}{rrrrr}
\toprule
Width & Selected & Residual std (\units) & Trend range (\units) &
Recovery within 1 sample\\
\midrule
17 & No & 0.578 & 8.768 & 98.00\%\\
25 & No & 1.012 & 8.037 & 98.92\%\\
33 & No & 1.372 & 8.076 & 99.33\%\\
41 & \textbf{Yes} & 1.572 & 8.280 & \textbf{99.85\%}\\
Raw & -- & -- & -- & 68.67\%\\
\bottomrule
\end{tabular}
\end{table}

Width 41 (2.05 seconds at 20\,Hz) was selected by the highest training-only oracle recovery
within one sample. It is not the minimum-median-error width: width 25's median absolute period
error is 0.122 samples versus 0.216 for width 41. The oracle diagnostic uses the known synthetic
period only for training-window preprocessing selection/diagnosis; it is not fed to the
forecaster and would not be available as such in ordinary deployment.

\fig{extracted_images/04-cell-18-output-05.png}
{Output 2.2: oracle period recovery within one sample for the raw context and four decomposition
widths. The selected width maximizes this criterion.}
{fig:oracle}

\subsection{Output 2.3: decomposition ablation}

\begin{table}[H]
\centering
\caption{Output 2.3: validation errors by population and product.}
\label{tab:decomp-ablation}
\scriptsize
\begin{tabular}{llrrr}
\toprule
Model & Population / product & MSE $(\units)^2$ & RMSE (\units) & MAE (\units)\\
\midrule
Raw Attention & Established / A & 0.150 & 0.388 & 0.293\\
Raw Attention & Established / B & 0.233 & 0.482 & 0.358\\
Raw Attention & Established / C & 0.175 & 0.418 & 0.306\\
Raw Attention & Held-out / A & 0.418 & 0.647 & 0.502\\
Raw Attention & Held-out / B & 0.336 & 0.580 & 0.439\\
Raw Attention & Held-out / C & 0.369 & 0.607 & 0.456\\
\midrule
Attention + decomposition & Established / A & 0.070 & 0.264 & 0.204\\
Attention + decomposition & Established / B & 0.112 & 0.335 & 0.249\\
Attention + decomposition & Established / C & 0.080 & 0.283 & 0.212\\
Attention + decomposition & Held-out / A & 0.129 & 0.359 & 0.275\\
Attention + decomposition & Held-out / B & 0.238 & 0.488 & 0.357\\
Attention + decomposition & Held-out / C & 0.189 & 0.434 & 0.330\\
\bottomrule
\end{tabular}
\end{table}

\fig{extracted_images/05-cell-19-output-02.png}
{Output 2.3: validation ablation showing the effect of decomposition on forecasts and errors.}
{fig:decomp-ablation}

On aggregate validation, decomposition lowers Attention RMSE from 0.431 to 0.296 on established
stations and from 0.612 to 0.430 on station 4, reductions of about 31\% and 30\%. It improves
each product/population row in the ablation. A moving average is sensible for the synthetic
world because the baseline varies smoothly on a longer time scale than the vibration cycle.
An abrupt baseline jump would violate this assumption: averaging would smear the transition and
leave transient structure in the residual, where it may be mistaken for repetition.

\section{Part 3: mixing by temporal delay}

\subsection{Response 3: delay aggregation and recurrence evidence}

For each selected delay $\tau_j$ and weight $\alpha_j$, the aggregation implements
\[
z_t=\sum_{j=1}^{K}\alpha_j v_{(t-\tau_j)\bmod L}.
\]
The implementation vectorizes selection with circular indices and shares each example's delays
and weights across heads and features. The immediate worked example returns
$[35,15,25,25]$ for values $[10,20,30,40]$, delays $[1,3]$, and weights $[0.75,0.25]$.
This verifies the required direction: $z_0=0.75(40)+0.25(20)=35$.

\begin{table}[H]
\centering
\caption{Output 3.1: supplied FFT delay scores equal direct circular correlation.}
\label{tab:delay-check}
\begin{tabular}{rrr}
\toprule
Delay & FFT score & Direct score\\
\midrule
0 & 4 & 4\\
1 & -4 & -4\\
2 & 4 & 4\\
3 & -4 & -4\\
\bottomrule
\end{tabular}
\end{table}

\subsection{Output 3.2: recurrence diagnostic}

\fig{extracted_images/06-cell-31-output-02.png}
{Output 3.2: raw and decomposed-signal circular recurrence for Station 2/Product A. The
diagnostic is computed from signal values, not learned query/key scores.}
{fig:recurrence}

For the plotted run, the actual period is 16.135 samples. Residual recurrence peaks at integer
delays 16 and 32, near $P$ and $2P$, with correlations 0.896 and 0.897; the next peak is near
48, approximately $3P$. The raw signal's slow drift creates broad, less period-specific
correlation. Removing the smooth component makes the repeating structure more visible and
supports a mixer organized by displacement. This is a signal-level diagnostic, not evidence
that the trained model's learned query/key scores peak at those delays.

Delay mixing would be a poor assumption for an isolated transient: it has no recurring
displacement, and circularly shifted copies can wrap or spread the event to unrelated positions.

\section{Part 4: compare designs and deploy}

\subsection{Response 4(a): hypothesis}

\textbf{Timing note.} As with Response 1A, this prediction was entered after saved experiment
outputs existed and is presented as a retrospective hypothesis, not as a preregistered prediction.
The hypothesis was that decomposition should help most when the slow component is prominent,
because it removes the ramp from the encoder and can expose recurring structure. Delay mixing
alone was expected to benefit less directly because it still scores projected features from a
raw representation.

\subsection{Response 4(b): comparison and horizon behavior}

\subsubsection{Output 4.1: 2-by-2 comparison}

\begin{table}[H]
\centering
\caption{Output 4.1: all four neural configurations on validation.}
\label{tab:four-models}
\small
\begin{tabularx}{\linewidth}{>{\raggedright\arraybackslash}X*{6}{>{\centering\arraybackslash}X}}
\toprule
Model & Est. RMSE & Est. MSE & Held-out RMSE & Held-out MSE & Parameters & Fit (s)\\
\midrule
Raw Attention & 0.431 & 0.186 & 0.612 & 0.374 & 368,368 & 760.032\\
Attention + decomposition & 0.296 & 0.087 & 0.430 & 0.185 & 373,024 & 897.249\\
Raw delay mixer & 0.221 & 0.049 & 0.340 & 0.116 & 368,370 & 899.230\\
Autoformer-inspired & \textbf{0.199} & \textbf{0.040} & \textbf{0.239} & \textbf{0.057} & 373,026 & 1008.758\\
\bottomrule
\end{tabularx}
\end{table}

\begin{table}[H]
\centering
\caption{Output 4.1: product-wise validation RMSE for the four neural configurations.}
\label{tab:four-products}
\small
\begin{tabular}{llrrr}
\toprule
Population & Model & A & B & C\\
\midrule
Established & Raw Attention & 0.388 & 0.482 & 0.418\\
Established & Attention + decomposition & 0.264 & 0.335 & 0.283\\
Established & Raw delay mixer & 0.211 & 0.211 & 0.239\\
Established & Autoformer-inspired & 0.226 & 0.188 & 0.182\\
\midrule
Held-out & Raw Attention & 0.647 & 0.580 & 0.607\\
Held-out & Attention + decomposition & 0.359 & 0.488 & 0.434\\
Held-out & Raw delay mixer & 0.415 & 0.281 & 0.309\\
Held-out & Autoformer-inspired & 0.270 & 0.221 & 0.224\\
\bottomrule
\end{tabular}
\end{table}

\fig{extracted_images/07-cell-35-output-02.png}
{Output 4.1: contrasting validation windows with the largest combined-model gain and loss. The
combined model is not best on every individual window.}
{fig:contrasting}

Both switches improve aggregate validation error. For established stations, decomposition lowers
MSE by 0.099 with Attention and 0.009 with delay mixing; replacing Attention with delay mixing
lowers it by 0.137 on raw inputs and 0.047 on decomposed inputs. For station 4, the corresponding
reductions are 0.189/0.059 and 0.258/0.128. Delay mixing alone has the larger aggregate gain, but
these scores do not isolate the drift-prominent contexts targeted by the hypothesis. Combined
gains (0.146 established, 0.317 held-out) are smaller than the sums of individual gains
(0.236/0.447). This is consistent with overlapping benefits in this run, not proof of a causal
mechanism.

Autoformer-inspired is the strongest neural model in aggregate, but not on every product/window.
Established Product A favors raw delay mixing (RMSE 0.211 vs. 0.226). In the largest-loss
contrast, raw delay mixing has RMSE 0.139 versus 0.358 for Autoformer-inspired.

\subsubsection{Output 4.2: horizon behavior}

\fig{extracted_images/08-cell-36-output-02.png}
{Output 4.2: RMSE across the 48-step forecast horizon for established and held-out populations.}
{fig:horizon}

\begin{table}[H]
\centering
\caption{Output 4.2: RMSE at forecast steps 1 and 48 (validation, \units).}
\label{tab:horizon-endpoints}
\small
\begin{tabular}{llrr}
\toprule
Population & Model & Step 1 & Step 48\\
\midrule
Established & Period-routed ridge & 0.112 & 0.228\\
Established & Raw Attention & 0.144 & 0.528\\
Established & Raw delay mixer & 0.121 & 0.346\\
Established & Attention + decomposition & 0.149 & 0.399\\
Established & Autoformer-inspired & 0.135 & 0.237\\
\midrule
Held-out & Period-routed ridge & 0.119 & 0.212\\
Held-out & Raw Attention & 0.176 & 0.761\\
Held-out & Raw delay mixer & 0.139 & 0.630\\
Held-out & Attention + decomposition & 0.147 & 0.589\\
Held-out & Autoformer-inspired & 0.147 & 0.299\\
\bottomrule
\end{tabular}
\end{table}

Errors generally rise with forecast horizon, though not monotonically. The increase is steepest
for Raw Attention and Raw delay mixer, particularly on the held-out station. Autoformer-inspired
and period-routed ridge remain comparatively stable toward step 48.

\subsection{Response 4(c): deployment choice and test results}

Model choice was made using validation evidence. Period-routed ridge is selected for both
populations because it has the lowest validation error, 59,904 parameters, and a recorded fitting
time of approximately 0.006 seconds, far below the neural fits. Sensor-specific ridge cannot
serve station 4 because it has no fitted model for that station.

\begin{table}[H]
\centering
\caption{Output 4.3: selected model validation and final test results.}
\label{tab:deployment}
\begin{tabular}{l l r r r r}
\toprule
Population & Model & Val. MSE & Val. RMSE & Test MSE & Test RMSE\\
\midrule
Established & Period-routed ridge & 0.028 & 0.166 & 0.026 & 0.161\\
Held-out & Period-routed ridge & 0.030 & 0.173 & 0.033 & 0.183\\
\bottomrule
\end{tabular}
\end{table}

The close validation/test RMSEs support generalization within this simulated setting. Test scores
are reported as the final evaluation of a validation-selected model, not used to select it.

\section{Conclusion}

The experiment illustrates the value of inductive bias. Period-routed ridge is substantially
better than the other listed models on both validation populations despite being much smaller and
faster to fit. Among neural variants, decomposition and delay mixing each improve aggregate
validation accuracy relative to Raw Attention, and their combination performs best overall.
However, the improvements are subadditive in this run and individual product/window exceptions
remain. These observations support the methods on this synthetic benchmark; they do not establish
universal superiority or real-world equipment-monitoring performance.

\section*{Reproducibility and AI-use disclosure}
\addcontentsline{toc}{section}{Reproducibility and AI-use disclosure}

The figures in this report are the eight saved PNG outputs extracted from the executed full-preset
notebook. Numerical tables reproduce its saved CSV outputs. No new training or notebook execution
was performed to prepare this report. The recorded run uses data seed 0 and model seed 0; results
are therefore a single-seed experiment.

Generative AI assistance (Copilot) was used to help inspect Task 1 requirements, compare notebook
outputs with those requirements, extract notebook figures, and draft/organize explanatory report
text. Representative prompts requested an implementation/results audit, comparison of hypotheses
with validation results, and a comprehensive LaTeX report. The AI output included draft summaries,
interpretations, tables, and report structure. The hypotheses were supplied by the student; the
report explicitly identifies them as retrospective because they were added after saved outputs
existed. The student revised the 1A ordering language and supplied/revised the 4A drift-focused
hypothesis; the AI draft was adjusted to reflect those edits and the saved validation values. The
AI did not run training or generate the experimental measurements. The student should review this
report, verify all claims, and disclose any further edits or AI use in accordance with course
policy.

\section*{References}
\addcontentsline{toc}{section}{References}

\begin{enumerate}[leftmargin=*]
  \item Wu, H. et al. (2021). ``Autoformer: Decomposition Transformers with Auto-Correlation
  for Long-Term Series Forecasting,'' Sections 3.1--3.2.
  \item Vaswani, A. et al. (2017). ``Attention Is All You Need.''
  \item Zeng, A. et al. (2023). ``Are Transformers Effective for Time Series Forecasting?''
  \item Hyndman, R. J. and Athanasopoulos, G. \textit{Forecasting: Principles and Practice},
  Section 5.10.
  \item AI651 Assignment 1, ``Forecasting Across Heterogeneous Sensors,'' Fall 2026; supplied
  assignment notebook and harness.
\end{enumerate}

\end{document}
